\documentclass[10pt,a4paper]{amsart}
\usepackage[margin=19mm]{geometry}
\usepackage{amsmath,amssymb,mathtools}
\usepackage[scr=rsfs,cal=euler]{mathalfa}
\usepackage{xfrac}
\usepackage[unicode,hidelinks,bookmarks=false]{hyperref}
\newcommand{\ie}{i.e.\ }
\newcommand{\cf}{cf.\ }
\newcommand{\resp}{respectively\ }
\newcommand{\Subsection}{\subsection}
\providecommand{\nobreakdash}{\nobreak\hbox{-}}
\setlength{\emergencystretch}{2em}
\multlinegap0pt
\usepackage{amssymb}
\usepackage{mathtools}
\usepackage[shortlabels]{enumitem}
\usepackage{tikz}
\newtheorem{theorem}{Theorem}
\newtheorem{restatable}{Restatable}
\newtheorem{claim}{Claim}
\newtheorem{corollary}{Corollary}
\newcommand{\NN}{\mathbb{N}}
\renewcommand{\ge}{\geqslant}
\renewcommand{\geq}{\ge}
\renewcommand{\le}{\leqslant}
\renewcommand{\leq}{\le}
\title{The Multicolour Size Ramsey Number of a Path}
\author{Csongor Beke, Anqi Li, Julian Sahasrabudhe}
\date{Source: arXiv:2511.16656v2 (CC BY 4.0)}
\begin{document}
\maketitle

\noindent\emph{Source and licence.} The Multicolour Size Ramsey Number of a Path. Version arXiv:2511.16656v2 (CC BY 4.0), distributed by arXiv under the Creative Commons Attribution 4.0 International licence, http://creativecommons.org/licenses/by/4.0/, as recorded in the arXiv licence metadata for this exact version and shown on the abstract page https://arxiv.org/abs/2511.16656v2. Full licence notice: \texttt{licenses/arxiv-2511.16656-CC-BY-4.0.txt}.\par\smallskip

\noindent\textit{Editorial scope.} Counted unit: the article's theorem thm:main (source0.tex lines 307--310). The article's complete original proof of it is delivered with it (source0.tex lines 522--623, one \texttt{proof} environment of 102 lines, which the article places away from the statement). No mathematics is written, repaired or re-derived: the statement and the proof are the article's own lines, in the article's own notation, and no retained line is substituted or re-flowed.  Retained uncounted as the source-local context the proof needs: lines 444--449.  Nothing was omitted from any retained span, so the package carries no omission record.  No retained line contains a citation, so the wrapper carries no bibliography and no bibliography entry.  No retained line contains a figure environment, an image-inclusion command or drawing code, so no image is delivered and the package declares no asset dependency.  Editorial formatting only, in addition: an AMS \texttt{amsart} preamble that restates the article's own declarations and macros used by the retained lines, namely the environments \texttt{theorem}, \texttt{restatable}, \texttt{claim}, \texttt{corollary} on separate counters, together with the standard packages those lines need.  The retained environments print in this document as Theorem 1, Restatable 1, Claim 1, Corollary 1 in order of appearance, whereas the article prints the same items with the numbering of its own full document.  The complete native source of the article as distributed in the arXiv e-print for this exact version is bundled unchanged as \texttt{sources/189526/source0.tex}.
\begin{restatable}[Key lemma]{lemma}{mainlemma}\label{lem:main}
There exists a constant $\beta_0 \in (0,1]$ such that for all sufficiently large $r$ and $k\geq 200\log r$ we have the following. Let $G$ be a graph with 
\[e(G) \leq (r^2 \log r) \, k\hspace{0.8em} \text{and} \hspace{0.8em}\Delta(G) = \beta r \log r\in \NN \]
where $\beta \in (0,\beta_0]$.
Then there exists a $P_k$-free subgraph $H \subset G$ such that \[e(H) \geq \dfrac{10 e(G)}{\beta^{0.9} r}.\] 
\end{restatable}

\begin{theorem}\label{thm:main}
    For any $r$ and $k \geq 200\log r$, we have 
    $$\widehat{R}_r(P_k)=\Theta\left((r^2 \log r) \,k\right).$$  
\end{theorem} 

\begin{proof}[Proof of Theorem~\ref{thm:main}] It suffices to prove the theorem for sufficiently large $r$, since the finitely many smaller values of $r$ can be absorbed into the choice of the constant $c_0$. Let $\beta_0$ be given by Lemma~\ref{lem:main}, let $r$ be sufficiently large  and take $\alpha=8/9$. We let
\[ c_0 \leq  \frac{1-\alpha}{10}\beta_0 \qquad \text{and} \qquad k \geq 200\log r,\]
and let $G$ be a graph with $e(G)\leq c_0(r^2\log r)k$ edges. We show that there is an $r$-colouring of $G$ with no monochromatic $P_k$.

First, we iteratively remove vertices of $G$ that have degree higher than $\beta_0r\log r$, and move them into a set $V_{-1}$. Note that
\[|V_{-1}|\leq\frac{c_0r^2(\log r)k}{\beta_0r\log r}\leq \frac{(1-\alpha)rk}{10}\leq \frac{rk}{90}.\]
Let the resulting graph be $G_0$. Then $G_0$ satisfies
\[e(G_0)\leq c_0\, r^2(\log r)k\qquad\text{and}\qquad
\Delta(G_0)\leq \beta_0\, r(\log r).\]

For each $i\geq 0$, we will construct a graph $G_i$ with parameters
\begin{equation}\label{eq:prop-Gi} e(G_i)=c_i\, r^2(\log r)k
\qquad\text{and}\qquad
\Delta(G_i)\leq \beta_i\, r(\log r), \qquad\text{where}\qquad \beta_i=\beta_0\Bigl(\frac{c_i}{c_0}\Bigr)^\alpha,\end{equation}
as long as $\beta_i\geq1/(r\log r)$. Once $\beta_i<1/(r\log r)$, we terminate the process and set $t=i$. Note that $\Delta(G_t)<1$, so $G_t$ is the empty graph.

Given $G_i$ satisfying \eqref{eq:prop-Gi}, we may apply Lemma~\ref{lem:main} to $G_i$ with $\beta=\beta_i$, obtaining a $P_k$-free subgraph $H_i\subset G_i$ with
\[e(H_i)\geq \frac{10\,e(G_i)}{\beta_i^{0.9}r}.\]

Set $F_{i,0}:=G_i\setminus H_i$. For $j\geq 0$, while $F_{i,j}$ contains a vertex of degree more than
\[\beta_0\Bigl(\frac{e(F_{i,j})}{c_0 r^2(\log r)k}\Bigr)^\alpha r(\log r),\]
we delete one such vertex and call the resulting graph $F_{i,j+1}$.
When this process stops, let $V_i$ be the set of deleted vertices, let $G_{i+1}=F_{i,|V_i|}$ and write $e(G_{i+1})=c_{i+1}\,r^2(\log r)k$. We
 then have by construction
\[\Delta(G_{i+1})\leq \beta_0\Bigl(\frac{c_{i+1}}{c_0}\Bigr)^\alpha r(\log r)
=\beta_{i+1}r(\log r).\]

In order to bound the number of iterations, we will need the following recursion on $\beta_i$.

\begin{claim}\label{claim:ci}
    We claim that 
    \[\beta_i^{0.9}\leq \beta_0^{0.9}-\frac{8 i}{r}.\]
\end{claim}

\begin{proof}
    Since $e(G_{i+1})\leq e(F_{i,0})=e(G_i)-e(H_i)$, we have
\[c_{i+1}\leq c_i\left(1-\frac{10}{\beta_i^{0.9}r}\right).\]
As long as the process has not terminated, we have $\beta_i\geq 1/(r\log r)$, and hence $\frac{10}{\beta_i^{0.9}r} <1$ for all sufficiently large $r$. Using $(1-x)^{0.8}\leq 1-0.8x$ for $x\in[0,1]$, we obtain
\[\beta_{i+1}^{0.9}=
\beta_0^{0.9}\Bigl(\frac{c_{i+1}}{c_0}\Bigr)^{0.8} \leq
\beta_0^{0.9}\Bigl(\frac{c_i}{c_0}\Bigr)^{0.8}
\left(1-\frac{10}{\beta_i^{0.9}r}\right)^{0.8} =
\beta_i^{0.9}\left(1-\frac{10}{\beta_i^{0.9}r}\right)^{0.8}\leq
\beta_i^{0.9}-\frac{8}{r}.\]
Iterating, we get
\[\beta_i^{0.9}\leq \beta_0^{0.9}-\frac{8 i}{r},\]
as desired.
\end{proof}

\begin{corollary}
    The iteration stops with $t \leq r/8$.
\end{corollary}
\begin{proof}
    As $\beta_t$ is a positive real number, and $\beta_0<1$, we get $t\leq r\cdot\beta_0^{0.9}/8\leq r/8$.
\end{proof}

Next, we control the total number of deleted vertices.

\begin{claim}
   We have
    \[\sum_{i=0}^{t-1}|V_i|\leq \frac{rk}{10}.\]
\end{claim}

\begin{proof} 
 For $0\leq j\leq |V_i|$, let $e(F_{i,j})=d_{i,j}\,r^2(\log r)k$, so when the vertex deleted in passing from $F_{i,j}$ to $F_{i,j+1}$ is removed, its degree is more than $\beta_0\Bigl(\frac{d_{i,j}}{c_0}\Bigr)^\alpha r(\log r)$. Hence
\[(d_{i,j}-d_{i,j+1})\,r^2(\log r)k
\geq \beta_0\Bigl(\frac{d_{i,j}}{c_0}\Bigr)^\alpha r(\log r),\]
and therefore
\[\frac{\beta_0}{rkc_0^{\alpha}}
\leq \frac{d_{i,j}-d_{i,j+1}}{d_{i,j}^\alpha}\leq \int_{d_{i,j+1}}^{d_{i,j}}x^{-\alpha}\,dx
=\frac{d_{i,j}^{1-\alpha}-d_{i,j+1}^{1-\alpha}}{1-\alpha},\]
where we used that $x\mapsto x^{-\alpha}$ is decreasing. Note that $d_{i,0}
\leq c_i$ and $d_{i,|V_i|}=c_{i+1}$.
Summing over $0\leq j\leq |V_i|$, we obtain
\[\frac{|V_i|}{rk}
\leq \frac{c_0^\alpha}{\beta_0(1-\alpha)}
\sum_{j=0}^{|V_i|-1}\bigl(d_{i,j}^{1-\alpha}-d_{i,j+1}^{1-\alpha}\bigr)
\leq \frac{c_0^\alpha}{\beta_0(1-\alpha)}
\bigl(c_i^{1-\alpha}-c_{i+1}^{1-\alpha}\bigr).\]
By summing over $0\leq i<t$, we conclude that
\[\sum_{i=0}^{t-1}\frac{|V_i|}{rk}
\leq
\frac{c_0^\alpha}{\beta_0(1-\alpha)}
\sum_{i=0}^{t-1}\bigl(c_i^{1-\alpha}-c_{i+1}^{1-\alpha}\bigr)
\leq
\frac{c_0}{\beta_0(1-\alpha)}
\leq \frac1{10},\]
establishing the second part of the claim.    
\end{proof}

Now, we are ready to describe the colouring of $G$. Note that the edges of $G$ can be partitioned as
\[E(G)=E(R)\cup\bigcup_{i=0}^{t-1}E(H_i),\]
where $R$ is a subgraph of $G$ which has a vertex cover given by $C=\bigcup_{i=-1}^{t-1}V_i$.


We first colour $\bigcup_{i=0}^{t-1}E(H_i)$. Since $t\leq r/8$ and each $H_i$ is $P_k$-free, assigning a distinct colour to each $H_i$ gives a monochromatic $P_k$-free colouring using at most $r/8$ colours.

It remains to colour $E(R)$. Note that \[|C|=\sum_{i=-1}^{t-1}|V_i|\leq \frac{rk}{90}+\frac{rk}{10}=\frac{rk}{9},\] we can partition $C$ into at most $r/4$ sets, each of size less than $k/2-1$. By assigning each edge of $R$ to a colour corresponding to the index of either one of its endpoints, we obtain a colouring of $E(R)$ with $r/4$ colours, where we break ties arbitrarily. Each colour class is a vertex-disjoint union of at most $k/2-1$ stars, so this gives a monochromatic $P_k$-free colouring.

Combining the two colourings above gives a monochromatic $P_k$-free colouring of $E(G)$
using at most $r/4+r/8<r$ colours
\end{proof}

\end{document}
